\documentclass[12pt,a4paper]{article}

\usepackage[utf8]{inputenc}
\usepackage[vietnamese]{babel}
\usepackage{amsmath,amssymb,amsthm}
\usepackage{geometry}
\usepackage{enumitem}
\usepackage[most]{tcolorbox}
\usepackage{hyperref}
\usepackage{fancyhdr}
\usepackage{tikz}
\usepackage{eso-pic}
\usepackage{xcolor}
\usepackage{array}
\usepackage{booktabs}

\geometry{a4paper, margin=2cm, bottom=2.5cm}
\hypersetup{colorlinks=true, linkcolor=blue!70!black, urlcolor=blue}
\setlist[itemize]{topsep=4pt, itemsep=2pt}
\setlist[enumerate]{topsep=4pt, itemsep=2pt}

\AddToShipoutPictureFG{%
  \AtPageCenter{%
    \begin{tikzpicture}[remember picture, overlay]
      \node[rotate=45, scale=1, opacity=0.06]{\fontsize{60}{72}\selectfont\bfseries\sffamily Đặng Tấn Quí};
    \end{tikzpicture}%
  }%
}

\pagestyle{fancy}
\fancyhf{}
\fancyhead[L]{\small\textit{Ma trận và Hệ phương trình}}
\fancyhead[R]{\small\textit{Đặng Tấn Quí}}
\fancyfoot[C]{\thepage}
\fancyfoot[L]{\footnotesize\textcopyright\ Đặng Tấn Quí}
\fancyfoot[R]{\footnotesize SĐT: 0377\,122\,210}
\renewcommand{\headrulewidth}{0.4pt}
\renewcommand{\footrulewidth}{0.4pt}
\fancypagestyle{plain}{\fancyhf{}\fancyfoot[C]{\thepage}\fancyfoot[L]{\footnotesize\textcopyright\ Đặng Tấn Quí}\fancyfoot[R]{\footnotesize SĐT: 0377\,122\,210}\renewcommand{\headrulewidth}{0pt}\renewcommand{\footrulewidth}{0.4pt}}

\newtcolorbox{dinhnghia}[1][]{enhanced, breakable, colback=blue!3!white, colframe=blue!60!black, fonttitle=\bfseries, title={Định nghĩa}, #1}
\newtcolorbox{dinhly}[1][]{enhanced, breakable, colback=green!3!white, colframe=green!50!black, fonttitle=\bfseries, title={Định lý}, #1}
\newtcolorbox{tinhchat}[1][]{enhanced, breakable, colback=yellow!5!white, colframe=orange!50!black, fonttitle=\bfseries, title={Tính chất}, #1}
\newtcolorbox{phuongphap}[1][]{enhanced, breakable, colback=gray!5!white, colframe=gray!50!black, fonttitle=\bfseries, title={Phương pháp}, #1}
\newtcolorbox{luuy}[1][]{enhanced, breakable, colback=red!3!white, colframe=red!50!black, fonttitle=\bfseries, title={Lưu ý}, #1}
\newtcolorbox{congthuc}[1][]{enhanced, breakable, colback=cyan!3!white, colframe=cyan!50!black, fonttitle=\bfseries, title={Công thức}, #1}
\newtcolorbox{vidu}[1][]{enhanced, breakable, colback=violet!3!white, colframe=violet!50!black, fonttitle=\bfseries, title={Ví dụ}, #1}

\begin{document}

\thispagestyle{empty}
\begin{tikzpicture}[remember picture, overlay]
  \fill[blue!80!black] (current page.south west) rectangle (current page.north east);
  \fill[blue!60!cyan, opacity=0.8] (current page.south west) -- (current page.north west) -- ([xshift=-3cm]current page.north east) -- (current page.south east) -- cycle;
  \foreach \r/\o in {6/0.05, 4.5/0.08, 3/0.12} {\fill[white, opacity=\o] ([xshift=2cm, yshift=-2cm]current page.north east) circle (\r cm);}
  \draw[white, line width=2pt] ([yshift=-5cm]current page.north west) ++(2cm,0) -- ++(17cm,0);
  \draw[white, line width=2pt] ([yshift=-16cm]current page.north west) ++(2cm,0) -- ++(17cm,0);
  \node[anchor=west, white, font=\fontsize{26}{32}\bfseries\selectfont] at ([xshift=3cm, yshift=-7.5cm]current page.north west) {PHẦN 8: MA TRẬN};
  \node[anchor=west, white, font=\fontsize{22}{28}\bfseries\selectfont] at ([xshift=3cm, yshift=-9.2cm]current page.north west) {VÀ HỆ PHƯƠNG TRÌNH};
  \node[anchor=west, white, font=\fontsize{16}{20}\selectfont] at ([xshift=3cm, yshift=-11cm]current page.north west) {Tài liệu luyện thi du học};
  \node[white, opacity=0.08, font=\fontsize{60}{60}\selectfont, rotate=-5] at ([xshift=-2cm, yshift=3cm]current page.center) {$\begin{bmatrix} a & b \\ c & d \end{bmatrix}$};
  \node[anchor=west, white, font=\Large\bfseries] at ([xshift=3cm, yshift=-13cm]current page.north west) {Biên soạn:};
  \node[anchor=west, yellow!80!white, font=\fontsize{28}{34}\bfseries\selectfont] at ([xshift=3cm, yshift=-14.5cm]current page.north west) {ĐẶNG TẤN QUÍ};
  \node[anchor=west, white!90, font=\large] at ([xshift=3cm, yshift=-18cm]current page.north west) {SĐT: 0377\,122\,210};
  \node[anchor=south, white!80, font=\small] at ([yshift=1.5cm]current page.south) {Tài liệu có bản quyền --- Cấm sao chép khi chưa được phép};
  \node[anchor=south east, white, font=\Large\bfseries] at ([xshift=-2cm, yshift=3cm]current page.south east) {\the\year};
\end{tikzpicture}
\newpage

\thispagestyle{empty}
\vspace*{5cm}
\begin{center}
  {\Large\bfseries PHẦN 8: MA TRẬN VÀ HỆ PHƯƠNG TRÌNH}\\[8pt]
  {\large Tài liệu luyện thi du học}
\end{center}
\vspace{2cm}
\begin{tcolorbox}[enhanced, colback=red!3!white, colframe=red!60!black, fonttitle=\bfseries, title={Thông tin bản quyền}, boxrule=1.5pt]
  \textbf{Tác giả:} Đặng Tấn Quí \quad \textbf{Liên hệ:} 0377\,122\,210 \quad \textbf{Năm:} \the\year \\[6pt]
  \textbf{Bản quyền \textcopyright\ \the\year\ Đặng Tấn Quí. Bảo lưu mọi quyền.}
\end{tcolorbox}
\vfill\begin{center}\small\textit{Tài liệu lưu hành nội bộ}\end{center}
\newpage
\tableofcontents
\newpage

%% ================================================================
%%  PHẦN 8: MA TRẬN VÀ HỆ PHƯƠNG TRÌNH
%% ================================================================
\section{Ma trận và Hệ phương trình}

%% ----------------------------------------------------------------
%%  8.1 Phép toán ma trận
%% ----------------------------------------------------------------
\subsection{Phép toán ma trận: Cộng, trừ, nhân ma trận, ma trận chuyển vị}

\begin{dinhnghia}[title={Ma trận}]
  \textbf{Ma trận $m \times n$} là bảng số hình chữ nhật gồm $m$ hàng và $n$ cột:
  \[
    A = \begin{pmatrix} a_{11} & a_{12} & \cdots & a_{1n} \\ a_{21} & a_{22} & \cdots & a_{2n} \\ \vdots & \vdots & \ddots & \vdots \\ a_{m1} & a_{m2} & \cdots & a_{mn} \end{pmatrix},
  \]
  ký hiệu $A = (a_{ij})$ hoặc $A \in M_{m\times n}$.

  \textbf{Các loại ma trận đặc biệt:}
  \begin{itemize}
    \item \textbf{Ma trận vuông:} $m = n$.
    \item \textbf{Ma trận đơn vị $I_n$:} đường chéo chính bằng 1, còn lại bằng 0.
    \item \textbf{Ma trận không $O$:} tất cả phần tử bằng 0.
    \item \textbf{Ma trận đường chéo:} tất cả phần tử ngoài đường chéo chính bằng 0.
  \end{itemize}
\end{dinhnghia}

\begin{congthuc}[title={Cộng và trừ ma trận}]
  Hai ma trận cùng cỡ ($m \times n$) có thể cộng/trừ theo từng phần tử:
  \[
    (A \pm B)_{ij} = a_{ij} \pm b_{ij}.
  \]
  Nhân với hằng số: $(kA)_{ij} = k a_{ij}$.
\end{congthuc}

\begin{congthuc}[title={Nhân ma trận}]
  Ma trận $A$ cỡ $m \times p$ và $B$ cỡ $p \times n$ có tích $AB$ cỡ $m \times n$:
  \[
    (AB)_{ij} = \sum_{k=1}^p a_{ik} b_{kj}.
  \]
  \textbf{Quy tắc:} Phần tử hàng $i$ cột $j$ của $AB$ = \textit{tích vô hướng} của hàng $i$ của $A$ với cột $j$ của $B$.
\end{congthuc}

\begin{luuy}
  \begin{itemize}
    \item Nhân ma trận \textbf{không giao hoán}: $AB \ne BA$ nói chung.
    \item $A(BC) = (AB)C$ (kết hợp); $A(B+C) = AB + AC$ (phân phối).
    \item $AI = IA = A$ với $I$ là ma trận đơn vị phù hợp.
    \item $AB = O$ không suy ra $A = O$ hoặc $B = O$.
  \end{itemize}
\end{luuy}

\begin{congthuc}[title={Ma trận chuyển vị}]
  \textbf{Chuyển vị} $A^T$ của $A = (a_{ij})$ cỡ $m \times n$ là ma trận cỡ $n \times m$ với $(A^T)_{ij} = a_{ji}$ (đổi hàng thành cột).

  Tính chất:
  \begin{itemize}
    \item $(A^T)^T = A$.
    \item $(A + B)^T = A^T + B^T$.
    \item $(AB)^T = B^T A^T$.
    \item $(kA)^T = k A^T$.
  \end{itemize}
\end{congthuc}

\begin{vidu}
  Cho $A = \begin{pmatrix}1 & 2 \\ 3 & 4\end{pmatrix}$ và $B = \begin{pmatrix}0 & 1 \\ -1 & 2\end{pmatrix}$. Tính $AB$, $BA$, và $A^T$.

  \medskip\textbf{Giải:}
  \[
    AB = \begin{pmatrix}1\cdot0+2\cdot(-1) & 1\cdot1+2\cdot2 \\ 3\cdot0+4\cdot(-1) & 3\cdot1+4\cdot2\end{pmatrix} = \begin{pmatrix}-2 & 5 \\ -4 & 11\end{pmatrix}.
  \]
  \[
    BA = \begin{pmatrix}0\cdot1+1\cdot3 & 0\cdot2+1\cdot4 \\ -1\cdot1+2\cdot3 & -1\cdot2+2\cdot4\end{pmatrix} = \begin{pmatrix}3 & 4 \\ 5 & 6\end{pmatrix}.
  \]
  $AB \ne BA$. \quad $A^T = \begin{pmatrix}1 & 3 \\ 2 & 4\end{pmatrix}$.
\end{vidu}

%% ----------------------------------------------------------------
%%  8.2 Ma trận nghịch đảo
%% ----------------------------------------------------------------
\subsection{Ma trận nghịch đảo: Cách tìm nghịch đảo 2x2 và 3x3}

\begin{dinhnghia}[title={Ma trận nghịch đảo}]
  Ma trận vuông $A$ cỡ $n \times n$ gọi là \textbf{khả nghịch} nếu tồn tại ma trận $A^{-1}$ sao cho:
  \[
    A A^{-1} = A^{-1} A = I_n.
  \]
  Ma trận không khả nghịch gọi là \textbf{suy biến (singular)}. Điều kiện: $\det(A) \ne 0$.
\end{dinhnghia}

\begin{congthuc}[title={Nghịch đảo ma trận $2\times 2$}]
  Cho $A = \begin{pmatrix}a & b \\ c & d\end{pmatrix}$, $\det(A) = ad - bc \ne 0$:
  \[
    A^{-1} = \frac{1}{ad-bc}\begin{pmatrix}d & -b \\ -c & a\end{pmatrix}.
  \]
  \textbf{Ghi nhớ:} Đổi chỗ $a$ và $d$, đổi dấu $b$ và $c$, chia cho $\det(A)$.
\end{congthuc}

\begin{phuongphap}[title={Nghịch đảo ma trận $3\times 3$ bằng phép khử Gauss-Jordan}]
  Dùng \textbf{ma trận bổ sung} $[A \mid I]$, thực hiện biến đổi hàng cơ bản (elementary row operations) đến khi vế trái thành $I$, lúc đó vế phải sẽ là $A^{-1}$:
  \[
    [A \mid I] \xrightarrow{\text{biến đổi hàng}} [I \mid A^{-1}].
  \]
  \textbf{Biến đổi hàng cơ bản:}
  \begin{itemize}
    \item Hoán đổi hai hàng: $R_i \leftrightarrow R_j$.
    \item Nhân hàng với hằng số $\ne 0$: $R_i \leftarrow k R_i$.
    \item Cộng bội số của một hàng vào hàng khác: $R_i \leftarrow R_i + k R_j$.
  \end{itemize}
\end{phuongphap}

\begin{congthuc}[title={Nghịch đảo bằng công thức định thức (ma trận phụ hợp)}]
  \[
    A^{-1} = \frac{1}{\det(A)} \text{adj}(A),
  \]
  trong đó $\text{adj}(A) = (C_{ij})^T$ là \textbf{ma trận phụ hợp} (matrix of cofactors transposed).

  Phần tử phụ đại số (cofactor): $C_{ij} = (-1)^{i+j} M_{ij}$, với $M_{ij}$ là định thức của ma trận con khi bỏ hàng $i$ và cột $j$.
\end{congthuc}

\begin{vidu}
  Tìm $A^{-1}$ với $A = \begin{pmatrix}1 & 2 & 0 \\ 0 & 1 & 1 \\ 1 & 0 & 1\end{pmatrix}$ bằng phép khử Gauss-Jordan.

  \medskip\textbf{Giải:}
  Lập $[A \mid I]$:
  \[
    \left[\begin{array}{ccc|ccc} 1&2&0&1&0&0 \\ 0&1&1&0&1&0 \\ 1&0&1&0&0&1 \end{array}\right]
    \xrightarrow{R_3 \leftarrow R_3-R_1}
    \left[\begin{array}{ccc|ccc} 1&2&0&1&0&0 \\ 0&1&1&0&1&0 \\ 0&-2&1&-1&0&1 \end{array}\right]
  \]
  \[
    \xrightarrow{R_3 \leftarrow R_3+2R_2}
    \left[\begin{array}{ccc|ccc} 1&2&0&1&0&0 \\ 0&1&1&0&1&0 \\ 0&0&3&-1&2&1 \end{array}\right]
    \xrightarrow{R_3 \leftarrow R_3/3}
    \left[\begin{array}{ccc|ccc} 1&2&0&1&0&0 \\ 0&1&1&0&1&0 \\ 0&0&1&-1/3&2/3&1/3 \end{array}\right]
  \]
  \[
    \xrightarrow{R_2 \leftarrow R_2-R_3}
    \left[\begin{array}{ccc|ccc} 1&2&0&1&0&0 \\ 0&1&0&1/3&1/3&-1/3 \\ 0&0&1&-1/3&2/3&1/3 \end{array}\right]
    \xrightarrow{R_1 \leftarrow R_1-2R_2}
    \left[\begin{array}{ccc|ccc} 1&0&0&1/3&-2/3&2/3 \\ 0&1&0&1/3&1/3&-1/3 \\ 0&0&1&-1/3&2/3&1/3 \end{array}\right]
  \]
  Vậy $A^{-1} = \dfrac{1}{3}\begin{pmatrix}1&-2&2\\1&1&-1\\-1&2&1\end{pmatrix}$.
\end{vidu}

%% ----------------------------------------------------------------
%%  8.3 Biến đổi hình học bằng ma trận
%% ----------------------------------------------------------------
\subsection{Biến đổi hình học: Quay, phản chiếu và co giãn bằng ma trận}

\begin{congthuc}[title={Ma trận biến đổi hình học trong $\mathbb{R}^2$}]
  Mỗi phép biến đổi tuyến tính $T: \mathbb{R}^2 \to \mathbb{R}^2$ được biểu diễn bởi một ma trận $2\times 2$.

  \textbf{Phép quay (Rotation) ngược chiều kim đồng hồ góc $\theta$:}
  \[
    R_\theta = \begin{pmatrix}\cos\theta & -\sin\theta \\ \sin\theta & \cos\theta\end{pmatrix}.
  \]

  \textbf{Phép phản chiếu (Reflection):}
  \begin{itemize}
    \item Qua trục $Ox$: $\begin{pmatrix}1 & 0 \\ 0 & -1\end{pmatrix}$.
    \item Qua trục $Oy$: $\begin{pmatrix}-1 & 0 \\ 0 & 1\end{pmatrix}$.
    \item Qua đường $y = x$: $\begin{pmatrix}0 & 1 \\ 1 & 0\end{pmatrix}$.
    \item Qua đường $y = -x$: $\begin{pmatrix}0 & -1 \\ -1 & 0\end{pmatrix}$.
    \item Qua gốc tọa độ: $\begin{pmatrix}-1 & 0 \\ 0 & -1\end{pmatrix}$.
  \end{itemize}

  \textbf{Phép co giãn (Scaling):}
  \[
    \begin{pmatrix}k_x & 0 \\ 0 & k_y\end{pmatrix} \quad \text{(co giãn theo cả hai trục)}.
  \]

  \textbf{Phép trượt (Shear):}
  \[
    \begin{pmatrix}1 & k \\ 0 & 1\end{pmatrix} \quad \text{(shear ngang)}.
  \]
\end{congthuc}

\begin{luuy}
  Áp dụng biến đổi cho điểm $(x, y)$: tính $\begin{pmatrix}x' \\ y'\end{pmatrix} = M \begin{pmatrix}x \\ y\end{pmatrix}$.

  Kết hợp nhiều biến đổi: nhân các ma trận (theo thứ tự ngược lại với thứ tự áp dụng).
\end{luuy}

\begin{vidu}
  Quay điểm $(1, 0)$ ngược chiều kim đồng hồ góc $90°$.

  \medskip\textbf{Giải:}
  $R_{90°} = \begin{pmatrix}0 & -1 \\ 1 & 0\end{pmatrix}$.

  $\begin{pmatrix}x' \\ y'\end{pmatrix} = \begin{pmatrix}0 & -1 \\ 1 & 0\end{pmatrix}\begin{pmatrix}1 \\ 0\end{pmatrix} = \begin{pmatrix}0 \\ 1\end{pmatrix}$.

  Điểm $(1,0)$ quay thành $(0,1)$. ✓
\end{vidu}

%% ----------------------------------------------------------------
%%  8.4 Giải hệ phương trình bằng ma trận
%% ----------------------------------------------------------------
\subsection{Giải hệ phương trình: Ma trận và định thức, biện luận số nghiệm}

\begin{dinhnghia}[title={Hệ phương trình tuyến tính dạng ma trận}]
  Hệ $n$ phương trình $n$ ẩn viết dưới dạng ma trận:
  \[
    A\vec{x} = \vec{b}, \quad A \in M_{n\times n},\quad \vec{x} = (x_1, \ldots, x_n)^T,\quad \vec{b} = (b_1, \ldots, b_n)^T.
  \]
\end{dinhnghia}

\begin{phuongphap}[title={Phương pháp 1: Khử Gauss (Gaussian Elimination)}]
  \begin{enumerate}
    \item Viết \textbf{ma trận bổ sung} $[A \mid \vec{b}]$.
    \item Thực hiện biến đổi hàng cơ bản để đưa về \textbf{dạng bậc thang hàng (Row Echelon Form)}.
    \item Giải ngược (back substitution) để tìm nghiệm.
  \end{enumerate}
\end{phuongphap}

\begin{phuongphap}[title={Phương pháp 2: Dùng ma trận nghịch đảo}]
  Nếu $\det(A) \ne 0$: nghiệm duy nhất $\vec{x} = A^{-1}\vec{b}$.
\end{phuongphap}

\begin{phuongphap}[title={Phương pháp 3: Quy tắc Cramer}]
  Nếu $\det(A) \ne 0$, nghiệm được tính bởi:
  \[
    x_i = \frac{\det(A_i)}{\det(A)},
  \]
  trong đó $A_i$ là ma trận $A$ với cột $i$ thay bằng $\vec{b}$.
\end{phuongphap}

\begin{dinhnghia}[title={Định thức (Determinant)}]
  \textbf{Ma trận $2 \times 2$:}
  \[
    \det\begin{pmatrix}a&b\\c&d\end{pmatrix} = ad - bc.
  \]
  \textbf{Ma trận $3 \times 3$ (khai triển theo hàng đầu):}
  \[
    \det(A) = a_{11}C_{11} + a_{12}C_{12} + a_{13}C_{13},
  \]
  hoặc theo \textbf{quy tắc Sarrus}:
  \[
    \det\begin{pmatrix}a&b&c\\d&e&f\\g&h&i\end{pmatrix} = aei + bfg + cdh - ceg - bdi - afh.
  \]
\end{dinhnghia}

\begin{congthuc}[title={Biện luận số nghiệm hệ $3 \times 3$ (ý nghĩa hình học)}]
  Hệ $3$ phương trình $3$ ẩn tương đương với $3$ mặt phẳng trong không gian. Số nghiệm phụ thuộc vào vị trí tương đối:
  \begin{itemize}
    \item $\det(A) \ne 0$: \textbf{nghiệm duy nhất} (3 mặt phẳng cắt nhau tại 1 điểm).
    \item $\det(A) = 0$ và hệ tương thích: \textbf{vô số nghiệm} (ít nhất 2 mặt phẳng trùng hoặc 3 mặt phẳng cắt theo 1 đường thẳng chung).
    \item $\det(A) = 0$ và hệ vô nghiệm: \textbf{vô nghiệm} (các mặt phẳng song song hoặc không có điểm chung).
  \end{itemize}
\end{congthuc}

\begin{vidu}
  Giải hệ $\begin{cases} x + 2y - z = 1 \\ 2x - y + z = 3 \\ x + y + 2z = 2 \end{cases}$ bằng khử Gauss.

  \medskip\textbf{Giải:}
  \[
    \left[\begin{array}{ccc|c}1&2&-1&1\\2&-1&1&3\\1&1&2&2\end{array}\right]
    \xrightarrow{R_2 \leftarrow R_2-2R_1,\, R_3 \leftarrow R_3-R_1}
    \left[\begin{array}{ccc|c}1&2&-1&1\\0&-5&3&1\\0&-1&3&1\end{array}\right]
  \]
  \[
    \xrightarrow{R_2 \leftarrow R_2/(-5)}
    \left[\begin{array}{ccc|c}1&2&-1&1\\0&1&-3/5&-1/5\\0&-1&3&1\end{array}\right]
    \xrightarrow{R_3 \leftarrow R_3+R_2}
    \left[\begin{array}{ccc|c}1&2&-1&1\\0&1&-3/5&-1/5\\0&0&12/5&4/5\end{array}\right]
  \]
  $\dfrac{12}{5}z = \dfrac{4}{5} \Rightarrow z = \dfrac{1}{3}$.
  $y - \dfrac{3}{5} \cdot \dfrac{1}{3} = -\dfrac{1}{5} \Rightarrow y = 0$.
  $x + 0 - \dfrac{1}{3} = 1 \Rightarrow x = \dfrac{4}{3}$.

  Nghiệm: $x = \dfrac{4}{3}$, $y = 0$, $z = \dfrac{1}{3}$.
\end{vidu}

\medskip\noindent\textbf{Các dạng bài tập tổng hợp:}
\begin{enumerate}[label=\textbf{Dạng \arabic*:}, leftmargin=*, align=left]
  \item Tính tổng, tích ma trận; tính chuyển vị.
  \item Tính định thức $2\times2$ và $3\times3$.
  \item Tìm ma trận nghịch đảo bằng Gauss-Jordan và bằng công thức.
  \item Giải hệ bằng khử Gauss, ma trận nghịch đảo, hoặc quy tắc Cramer.
  \item Biện luận số nghiệm hệ phương trình theo tham số.
  \item Tìm biến đổi hình học tương ứng với ma trận cho trước và ngược lại.
\end{enumerate}

\end{document}
